\section{Background and Experimental Setup}
\label{sec:background}

\subsection{The R-Matrix Mechanism}

The R-matrix (ancestor closure matrix) is a binary $C \times C$ matrix
where $R_{ij} = 1$ iff class $i$ is an ancestor of class $j$ (including
$i = j$). It is derived from the transitive closure of the hierarchy
adjacency matrix $A$:

\begin{equation}
R_{ij} = \begin{cases}
1 & \text{if } i = j \text{ or } \exists \text{ path } i \rightarrow \ldots \rightarrow j \text{ in } A \\
0 & \text{otherwise}
\end{cases}
\label{eq:rdef}
\end{equation}

The R-matrix operates at two distinct points in the HMC pipeline:

\textbf{Mode 1: Training-time consistency loss (MC-loss).}
Following~\cite{giunchiglia2018hmc}, the R-matrix is applied during
training to constrain model outputs before computing the binary
cross-entropy loss. For each batch, the raw logits $\mathbf{o} \in
\mathbb{R}^{B \times C}$ are constrained via:

\begin{equation}
\mathbf{o}^{\text{constr}}_i = \max_{j: R_{ij}=1} \mathbf{o}_j
\label{eq:constr}
\end{equation}

which ensures that the constrained output for class $i$ is the maximum
score among itself and all its descendants. The BCE loss is then
computed on a combination of constrained outputs for positive labels and
raw outputs for negative labels.

\textbf{Mode 2: Inference-time reconciliation.}
At prediction time, the same bottom-up max-propagation is applied to the
model's raw scores:

\begin{equation}
p_i^{\text{rec}} = \max\left(p_i,\; \max_{j \in \text{children}(i)} p_j^{\text{rec}}\right)
\label{eq:reconcile}
\end{equation}

Processing levels from deepest to shallowest guarantees that $\forall
(i,j): \text{ancestor}(i,j) \implies p_i^{\text{rec}} \geq
p_j^{\text{rec}}$ with zero post-inference violations.

\subsection{Datasets}

We evaluate on three representative datasets covering distinct hierarchy
types:

\begin{table}[H]
\centering
\caption{Datasets used in the ablation study.}
\label{tab:datasets}
\begin{tabular}{lrrrrr}
\toprule
\textbf{Dataset} & \textbf{Hierarchy} & \textbf{Train} & \textbf{Test} &
\textbf{Classes} & \textbf{Depth} \\
\midrule
ArXiv & Shallow tree & 40,982 & 10,245 & 157 & 2 \\
cellcycle\_FUN & Deep tree & 2,476 & 1,281 & 500 & 4+ \\
cellcycle\_GO & DAG & 2,474 & 1,280 & 4,126 & 13 \\
\bottomrule
\end{tabular}
\end{table}

For the scaling laws analysis, we extend to all 25 built-in datasets
spanning text (ArXiv, WOS, AAPD), tabular trees (10 FunCat datasets),
Gene Ontology DAGs (9 datasets), and cross-domain tabular (Enron,
Diatoms, ImCLEF). Full dataset statistics are in Appendix~\ref{sec:appendix}.

\subsection{Experimental Configuration}

All experiments use the HMC-Torch platform~\cite{sette2026hmctorch} with
a 3-layer MLP (hidden dimension 512 for text, 256 for tabular), AdamW
optimizer ($\beta_1=0.9, \beta_2=0.999$, weight decay $10^{-5}$,
learning rate $10^{-4}$), and prevalence-weighted BCE loss. Text
datasets use frozen SPECTER2-base embeddings (768 dimensions). Tabular
datasets use standard-scaled features with mean imputation.

For the DAG dataset (cellcycle\_GO), we use the Block-Diagonal R-Matrix
sparse approximation~\cite{sette2026hmctorch} which achieves zero
hierarchy violations in $O(C+E)$ memory without materializing the dense
$O(C^2)$ matrix. All experiments use seed 42 unless otherwise noted, on
an NVIDIA RTX 3060 (12GB).

Metrics are micro-averaged F1 (with threshold sweep over $[0.1, 0.9]$
step 0.05) and micro-averaged AUPRC, computed only over evaluable
classes (excluding root nodes).
